% === Revised Section 4: Kolmogorov complexity ===

% NEW PREAMBLE ADDITIONS NEEDED:
% \usepackage[most]{tcolorbox}  % already loaded in the main file
% \tcbuselibrary{breakable,skins}  % already loaded in the main file
%
% \newtcolorbox{activity}[1][]{
%   breakable,
%   enhanced,
%   colback=blue!5,
%   colframe=blue!55!black,
%   coltitle=white,
%   fonttitle=\bfseries,
%   title={Activity},
%   boxrule=0.6pt,
%   left=6pt,right=6pt,top=4pt,bottom=4pt,
%   #1
% }

\section{Kolmogorov complexity}
\label{sec:kolmogorov}

\begin{intuition}
\textbf{One-line answer.}  The Kolmogorov complexity $K(x)$ of an
object $x$ is the length, in bits, of the shortest computer program
that prints $x$ and stops (\cref{def:kolmogorov}).  It is the
vendor-independent floor on how far $x$ can ever be compressed.

\textbf{Twin USB sticks.}  Imagine two 8\,GB USB sticks on your desk.
One is full of family JPGs; the other is full of bytes drawn from
\texttt{/dev/urandom}.  Both weigh the same; both read the same number
of bits off the bus.  Zip the first and you get roughly
$6$--$7\,\text{GB}$; zip the second and you get back almost exactly
$8\,\text{GB}$ --- sometimes \emph{larger}, once zip's header is added.
The difference is not in the hardware.  It is a property of the
\emph{contents}: the photos have structure a short program can
exploit, the random bytes do not.  Kolmogorov complexity is the
mathematical name for that difference.

\textbf{The recipe-over-the-phone picture.}  You are a chef.  A
colleague in another city must reproduce your dish exactly.  You
dictate the shortest instructions that do the job.  ``Boil water''
takes two words.  A souffl\'e takes a page.  A pattern of $10{,}000$
genuinely random sprinkles on a cake has no shorter description than
itself: you read out every position.  Such an object is
\emph{Kolmogorov random}; it has no exploitable structure.  Everything
else --- the digits of $\pi$, the works of Shakespeare, your customer
database --- has a recipe shorter than the raw data.

\textbf{A tiny worked example.}  The string \texttt{01010101\ldots} (a
million characters long) has an $\approx 30$-character recipe:
\texttt{print "01" 500000 times}.  A million genuine coin flips has no
such shortcut; its recipe is, to within a small constant, a million
characters.  Same disk size, wildly different Kolmogorov complexity.
\emph{Complexity measures structure, not size.}

\textbf{Invariance and uncomputability, previewed.}  Two different
programming languages disagree about the length of a recipe by at most
the size of a translator between them (\cref{thm:invariance}).  So
$K(x)$ is a property of the object, not your toolkit.  But no
algorithm, given $x$, can hand you its shortest program
(\cref{thm:uncomputable}); if one existed, it would decide the halting
problem.  $K(x)$ is a perfectly well-defined number that you can bound
from above --- by running zip, by training a model --- but never
compute exactly.

\textbf{Why a decision-maker should care.}  Your company's data has a
\emph{true}, vendor-independent information content.  No one can quote
it to you exactly; anyone who claims to has sold you the Brooklyn
Bridge.  But every real compressor --- zip, JPEG, your LLM --- gives
you an upper bound, and upper bounds are enough to tell snake oil from
signal.
\end{intuition}

\begin{intuition}[title={The compression sniff test}]
If a vendor claims to losslessly compress genuinely random or properly
encrypted data, they are selling snake oil --- Shannon closed that
door in 1948.  A truly random source has entropy equal to its raw
bit-length; no code can do better on average
(\cref{thm:source-coding}).  ``We compress encrypted traffic by
40\,\%'' is either marketing about a non-random \emph{side channel}
(timestamps, packet sizes, TLS record framing) --- which may be
legitimate --- or a claim that should not appear in a contract you
sign.  Activity~\ref{act:bingo} gives you a bingo card for the next
vendor demo.
\end{intuition}

\begin{activity}[title={Activity 1 --- Shortest-recipe game}]
\label{act:recipe}
\textbf{Setup (10 minutes, pairs).}  Each pair receives the four
strings below.  Write, in the pseudocode of your choice, the shortest
program you can that prints each exactly.  Score: total characters
summed across the four programs.  Lowest score wins.

\begin{enumerate}[label=(\alph*),leftmargin=*]
  \item \texttt{01} repeated $500$ times (total length $1000$).
  \item The first $100$ decimal digits of $\pi$.
  \item The filename \texttt{IMG\_0423.JPG}.
  \item $100$ bytes freshly drawn from \texttt{/dev/urandom},
        printed as a $200$-character hex string (distributed on paper).
\end{enumerate}

\textbf{Debrief.}  Strings (a), (b), (c) collapse to tiny programs: a
loop, a formula, and a literal constant shorter than the output.
String (d) does not collapse: the shortest program anyone in the room
produces is essentially \texttt{print "<the 200 hex chars>"}.  You
have just re-derived, by hand, the distinction between
Kolmogorov-compressible and Kolmogorov-random objects.
\end{activity}

Shannon's framework is probabilistic: entropy is a property of a
distribution, not of an individual string.  Kolmogorov's move
\citep{kolmogorov1965three} was to attach a complexity measure to a
\emph{single} object by asking for the length of its shortest
description in a fixed universal language.

Fix once and for all a \textbf{universal Turing machine} $U$ --- a
machine that, given as input a description $\langle M, w\rangle$ of
any Turing machine $M$ and input $w$, simulates $M$ on $w$.

\begin{definition}[Kolmogorov complexity]
\label{def:kolmogorov}
The \textbf{Kolmogorov complexity} of a finite binary string
$x \in \{0,1\}^*$, relative to $U$, is
\begin{equation}
  K_U(x)
  \;=\; \min\bigl\{\, |\pi| \,:\, \pi \in \{0,1\}^*,\ U(\pi) = x \,\bigr\},
  \label{eq:kolmogorov}
\end{equation}
where $|\pi|$ is the bit-length of the program $\pi$ and $U(\pi) = x$
means $U$ on input $\pi$ halts with output $x$.
\end{definition}

\begin{activity}[title={Activity 2 --- The gzip-vs-$K$ scoreboard}]
\label{act:scoreboard}
\textbf{Setup (5 minutes, demo or take-home).}  For each of the four
strings from Activity~\ref{act:recipe}, measure three numbers:

\begin{enumerate}[label=(\roman*),leftmargin=*]
  \item \textbf{Raw size} --- the bit-length of the string itself.
  \item \textbf{\texttt{gzip -9} size} --- bit-length of the
        compressed file (subtract the $\sim\!80$-bit gzip header for
        tiny inputs, or just note it).
  \item \textbf{Shortest-recipe bound} --- the bit-length of the
        winning program from Activity~\ref{act:recipe}, converted to
        bits at $8$\,bit/character.
\end{enumerate}

Plot the three as a bar chart per string.  You will see: for (a)--(c)
the recipe bound sits far below gzip; for (d) all three bars are
essentially the same height.  The ``recipe bound'' is itself only an
upper bound on $K(x)$ --- the true $K(x)$ is uncomputable
(\cref{thm:uncomputable}) --- but seeing gzip dominated by a
one-line Python program is the fastest way to feel in your bones that
$K$ lives far below any off-the-shelf compressor.
\end{activity}

\begin{theorem}[Invariance theorem]
\label{thm:invariance}
If $U_1$ and $U_2$ are two universal Turing machines, then there
exists a constant $c_{U_1,U_2}$\footnote{The constant $c_{U_1,U_2}$
is finite and does not grow with $|x|$.  Concretely: the smallest
self-contained CPython interpreter in C is a few hundred kilobytes of
source, and a minimal C compiler written in Python is comparable; so
for the pair (Python, C), $c_{U_1,U_2}$ is on the order of
$10^6$--$10^7$ bits.  For a megabyte-scale corpus ($\sim\!10^7$ bits)
this is already a non-trivial fraction; for a one-gigabyte ERP export
($\sim\!10^{10}$ bits) it is numerically negligible, and
``machine-independent up to an additive constant'' is a property you
can plan around.  For objects \emph{smaller} than the constant --- a
one-sentence log entry, a password --- it is not, and quoting ``the''
Kolmogorov complexity of such an object is meaningless without fixing
$U$.}, depending on $U_1$ and $U_2$ but not on $x$, such that
\begin{equation}
  \bigl| K_{U_1}(x) - K_{U_2}(x) \bigr| \;\le\; c_{U_1,U_2}
  \quad \text{for every } x \in \{0,1\}^*.
  \label{eq:invariance}
\end{equation}
\end{theorem}

\begin{proof}[Proof sketch]
Because $U_1$ is universal, there is a fixed-length binary program
$s_{2\to 1}$ --- an interpreter for $U_2$ written for $U_1$ --- that,
when prepended to any $U_2$-program $\pi$, causes $U_1$ to simulate
$U_2(\pi)$.  Thus every $U_2$-program $\pi$ of length $|\pi|$ yields a
$U_1$-program of length $|s_{2\to 1}| + |\pi|$ with the same output,
giving $K_{U_1}(x) \le K_{U_2}(x) + |s_{2\to 1}|$.  Symmetrically with
$s_{1\to 2}$.  Taking $c_{U_1,U_2} = \max(|s_{1\to 2}|, |s_{2\to 1}|)$
gives~\eqref{eq:invariance}.
\end{proof}

By convention we henceforth write $K(x)$ and treat the universal
simulator as the additive constant.

\begin{intuition}[title={The Berry paradox, dramatised}]
Consider the English phrase
\begin{quote}
``the smallest positive integer not nameable in fewer than twenty
English words.''
\end{quote}
That phrase contains nineteen words.  So if the integer it picks out
exists, we have just named it in nineteen words, contradicting the
very property we used to pick it.  The paradox is resolved by noting
that ``nameable in English'' is too slippery to be a mathematical
predicate.  But if we replace ``nameable in English'' with ``printed
by a program of length $\le m$'', the slipperiness vanishes --- and
the same self-referential trick becomes a clean proof that $K$ cannot
be computed.  That is the content of \cref{thm:uncomputable}.
\end{intuition}

\begin{theorem}[Uncomputability of $K$]
\label{thm:uncomputable}
There is no total computable function $f : \{0,1\}^* \to \mathbb{N}$
such that $f(x) = K(x)$ for all $x$.
\end{theorem}

\begin{proof}[Proof (Berry-paradox style)]
Suppose for contradiction such an $f$ exists.  Define the program
$P_m$, parameterised by an integer $m \in \mathbb{N}$, that does the
following: enumerate $\{0,1\}^*$ in shortlex order and, for each $x$,
compute $f(x)$ and output the first $x$ such that $f(x) > m$; then
halt.  Strings of complexity greater than $m$ exist (there are only
$2^m - 1$ programs of length less than $m$, but infinitely many
strings), so $P_m$ halts and outputs some $x_m^\star$ with
$K(x_m^\star) > m$.

The program $P_m$ has a fixed body of bit-length $B$ (independent of
$m$), plus a binary encoding of $m$ of at most
$\lceil \log_2(m+1) \rceil + O(1)$ bits (e.g.\ Elias delta coding).
So $|P_m| \le B + \log_2 m + O(1)$, giving
\begin{equation}
  K(x_m^\star) \;\le\; |P_m| \;\le\; B + \log_2 m + O(1).
  \label{eq:berry-upper}
\end{equation}
Combined with $K(x_m^\star) > m$,
\begin{equation}
  m \;<\; B + \log_2 m + O(1),
  \label{eq:berry-contra}
\end{equation}
which fails for every sufficiently large $m$ (since $m - \log_2 m
\to \infty$) --- a contradiction.
\end{proof}

\begin{activity}[title={Activity 3 --- Snake-oil bingo}]
\label{act:bingo}
\textbf{Setup (5 minutes, individual then plenary).}  Below is a
$3\times 3$ card of claims collected from real compression- and
AI-vendor pitches.  Mark each cell \textbf{S} (violates Shannon /
Kolmogorov --- snake oil), \textbf{L} (legitimate, exploits a
non-random side channel), or \textbf{?} (need to ask the vendor a
clarifying question).

\begin{center}
\small
\begin{tabular}{|p{0.3\linewidth}|p{0.3\linewidth}|p{0.3\linewidth}|}
\hline
``Lossless 10:1 on AES-256 ciphertext.'' &
``40\,\% reduction on TLS traffic via header deduplication.'' &
``Patented sub-Shannon entropy coding.'' \\
\hline
``AI-powered compression of already-compressed JPEGs.'' &
``Adaptive model shrinks your logs by 8$\times$.'' &
``Zero-loss reduction of random-number tables.'' \\
\hline
``Our LLM produces an $n$-bit summary of an $n^2$-bit corpus.'' &
``Universal compressor --- works on any input.'' &
``Dictionary coding of repetitive Kafka topics, typical 70\,\%.'' \\
\hline
\end{tabular}
\end{center}

\textbf{Debrief.}  Cells about ciphertext, already-compressed JPEGs,
random tables, and ``universal'' compressors are S.  Header
deduplication on TLS and dictionary coding on repetitive topics are L.
``AI-powered'' and ``sub-Shannon'' are ? until the vendor tells you
what distribution they are assuming.  The LLM summary cell is L
\emph{only} if lossy.
\end{activity}

\medskip

The upshot of \cref{thm:uncomputable} is operational: every practical
compressor --- gzip, PNG, FLAC, and indeed every neural language model
--- produces an \emph{upper bound} on $K(x)$, never its exact value.
What a decision-maker does with this on Monday morning: (i) treat
every compression ratio a vendor quotes as a claim about \emph{your}
data's distribution, not about the compressor's cleverness; (ii)
demand that ratios be reported on a held-out sample from your own
corpus, not the vendor's demo set; (iii) recognise that the gap
between your best available compressor and the true $K$ is where
better models --- and, in \cref{sec:solomonoff}, better
\emph{predictors} --- still have room to win.
